%% =============================================================================
%% Inferential Statistics Cheat Sheet
%% BCB744 — Quantitative Biology
%% Chapters 1--26
%% =============================================================================
\documentclass[9pt,landscape]{article}

\usepackage[a3paper,top=10mm,bottom=10mm,left=10mm,right=10mm]{geometry}
\usepackage[T1]{fontenc}
\usepackage[utf8]{inputenc}
\usepackage[osf]{ebgaramond}
\usepackage[sfdefault,scaled=0.95]{inter}
\usepackage[scaled=0.88]{FiraMono}
\usepackage{sfmath}
\usepackage[protrusion=true,expansion=true,tracking=true]{microtype}
\usepackage[table]{xcolor}
\usepackage{pagecolor}
\usepackage{array}
\usepackage{booktabs}
\usepackage{longtable}
\usepackage{multicol}
\usepackage{enumitem}
\usepackage{textcomp}
\usepackage{xstring}
\usepackage[most]{tcolorbox}

\definecolor{MyBackground}{HTML}{FFFDF0}
\definecolor{headbg}{RGB}{28,40,52}
\definecolor{headfg}{RGB}{255,255,255}
\definecolor{groupbg}{RGB}{230,220,200}
\definecolor{diagbg}{RGB}{255,250,230}
\definecolor{flowbg}{RGB}{235,245,240}
\definecolor{warnbg}{RGB}{252,236,233}
\definecolor{reportbg}{RGB}{233,242,250}
\definecolor{rulelight}{RGB}{180,188,198}
\definecolor{codeblue}{RGB}{0,80,160}
\definecolor{warnred}{RGB}{180,40,30}
\definecolor{notegrey}{RGB}{95,95,105}

\newcommand{\Rc}[1]{{\ttfamily\footnotesize\color{codeblue}#1}}
\newcommand{\Note}[1]{{\scriptsize\itshape\color{notegrey}#1}}
\newcommand{\CueIcon}[1]{%
  \IfStrEqCase{#1}{%
    {G}{$\mu$}%
    {Rk}{$\diamond$}%
    {Asc}{$\leftrightarrow$}%
    {Bin}{\texttt{0/1}}%
    {Cnt}{\#}%
    {Grp}{$\bullet\!\bullet$}%
    {Sm}{$\sim$}%
    {NL}{$f(x)$}%
    {Q}{$Q$}%
    {Pred}{$\hat{y}$}%
    {Any}{$\Delta$}%
  }[\textbf{?}]%
}
\newcommand{\Cue}[1]{\shortstack[c]{\CueIcon{#1}\\[-1pt]\scriptsize\textbf{#1}}}
\newlength{\tableblockwidth}
\setlength{\tableblockwidth}{370mm}
\newlength{\tableoffset}
\setlength{\tableoffset}{15mm}
\newlength{\headerleftwidth}
\setlength{\headerleftwidth}{200mm}
\newlength{\headergap}
\setlength{\headergap}{20mm}
\newlength{\headerrightwidth}
\setlength{\headerrightwidth}{150mm}
\newtcolorbox{decisionbox}{
  enhanced,
  boxrule=1pt,
  arc=2mm,
  colback=groupbg,
  colframe=headbg,
  title=\textbf{QUICK DECISION FLOW},
  fonttitle=\sffamily\bfseries\small,
  left=4mm,right=4mm,top=2mm,bottom=2mm
}
\newtcolorbox{infopanelbox}[2]{
  enhanced,
  width=\linewidth,
  boxrule=1pt,
  arc=2mm,
  colback=#1,
  colframe=headbg,
  colbacktitle=headbg,
  left=3mm,right=3mm,top=2mm,bottom=2mm,
  title={#2},
  fonttitle=\sffamily\bfseries\small,
  coltitle=headfg
}
\newtcolorbox{warnpanelbox}[1]{
  enhanced,
  width=\linewidth,
  boxrule=1pt,
  arc=2mm,
  colback=MyBackground,
  colframe=warnred,
  colbacktitle=warnred,
  coltitle=headfg,
  coltext=warnred,
  left=3mm,right=3mm,top=2mm,bottom=2mm,
  title={#1},
  fonttitle=\sffamily\bfseries\small
}

\setlength{\LTpre}{0pt}
\setlength{\LTpost}{0pt}
\setlength{\extrarowheight}{1pt}
\renewcommand{\arraystretch}{1.24}
\frenchspacing
\pagecolor{MyBackground}
\pagestyle{empty}

\begin{document}

\noindent\hspace*{\tableoffset}%
\begin{minipage}{\tableblockwidth}
\begin{minipage}[t]{\headerleftwidth}
  \vspace{0pt}
  \raggedright
  {\Huge\sffamily\bfseries\color{headbg} Inferential Statistics Cheat Sheet}\\
  {\large\itshape BCB744 --- Biostatistics Version 3 (3 April 2026)}\\
  {\small Professor A.J. Smit}\\
  {\small Chapters 1--5 and 26 establish foundations, workflow, and interpretation; the sheet below covers the inferential, diagnostic, and model-based procedures introduced across Chapters 6--25.}
\end{minipage}
\hspace*{\headergap}%
\begin{minipage}[t]{\headerrightwidth}
  \vspace*{-2.5mm}
  \begin{decisionbox}
    \small\sffamily
    \textbf{1. Response?} Continuous \(\rightarrow\) Gaussian methods; binary or proportion \(\rightarrow\) binomial GLM; count \(\rightarrow\) Poisson or quasi-Poisson.\\
    \textbf{2. Independent?} No \(\rightarrow\) paired procedure or mixed/grouped model.\\
    \textbf{3. Linear enough?} No \(\rightarrow\) polynomial, GAM, nonlinear, or quantile model as appropriate.
  \end{decisionbox}
\end{minipage}
\end{minipage}

\vspace{8pt}

{\sffamily\small
\begin{longtable}{>{\raggedright\arraybackslash}p{9mm} >{\centering\arraybackslash}p{14mm} >{\raggedright\arraybackslash}p{31mm} >{\raggedright\arraybackslash}p{49mm} >{\raggedright\arraybackslash}p{57mm} >{\raggedright\arraybackslash}p{68mm} >{\raggedright\arraybackslash}p{113mm}}
\rowcolor{headbg}
\color{headfg}\textbf{Ch.} &
\color{headfg}\textbf{Cue} &
\color{headfg}\textbf{Method} &
\color{headfg}\textbf{Use When} &
\color{headfg}\textbf{Null / Model Target} &
\color{headfg}\textbf{Main R Function(s)} &
\color{headfg}\textbf{Notes, Assumptions, or Follow-up} \\
\endfirsthead

\rowcolor{headbg}
\color{headfg}\textbf{Ch.} &
\color{headfg}\textbf{Cue} &
\color{headfg}\textbf{Method} &
\color{headfg}\textbf{Use When} &
\color{headfg}\textbf{Null / Model Target} &
\color{headfg}\textbf{Main R Function(s)} &
\color{headfg}\textbf{Notes, Assumptions, or Follow-up} \\
\endhead

\midrule
\multicolumn{7}{r}{\scriptsize Continued on next page} \\
\endfoot

\bottomrule
\endlastfoot

\rowcolor{headbg}\multicolumn{7}{l}{\color{headfg}\textbf{TESTING DIFFERENCES AND BASIC CHECKS}} \\
\rowcolor{groupbg}\multicolumn{7}{l}{\textbf{Diagnostics and assumption checks}} \\
\rowcolor{diagbg}6 & \Cue{G} &
Shapiro--Wilk normality test &
One sample, paired differences, or within-group values need a formal normality check before a Gaussian mean-based test. &
$H_0$: the checked values are compatible with a normal distribution. &
\Rc{shapiro.test(x)} &
Check values on the analysis scale. Read with histograms and Q--Q plots, not alone. \\

\rowcolor{diagbg}6 & \Cue{G} &
Levene's test for equal variance &
Two or more groups are being compared and homogeneity of spread needs checking. &
$H_0$: group variances are equal. &
\Rc{car::leveneTest(y \textasciitilde{} g, data = df)} &
Preferred over Bartlett when normality is doubtful. Most relevant before Student's \textit{t}-test and classical ANOVA. \\

\rowcolor{diagbg}6 & \Cue{G} &
Bartlett's test &
Several groups are compared and a normal-theory variance test is acceptable. &
$H_0$: group variances are equal. &
\Rc{bartlett.test(y \textasciitilde{} g, data = df)} &
More sensitive to non-normality than Levene. Use only when normality is already plausible. \\

\rowcolor{flowbg}6 & \Cue{G} &
Transformation of the response &
Assumptions fail on the raw scale but a transformed scale may stabilise spread or reduce skew. &
No null; the aim is a more defensible analysis scale. &
\Rc{mutate(log\_y = log(y))}\newline
\Rc{mutate(sqrt\_y = sqrt(y))} &
Transform before the final Gaussian analysis, then re-check assumptions. \\

\rowcolor{groupbg}\multicolumn{7}{l}{\textbf{Mean comparisons: \textit{t}-tests and rank analogues}} \\
7 & \Cue{G} &
One-sample \textit{t}-test &
One continuous sample is compared with a biologically meaningful reference mean. &
$H_0$: $\mu = \mu_0$. &
\Rc{t.test(x, mu = m0)} &
Directional alternatives use \Rc{alternative = "less"} or \Rc{"greater"}. \\

7 & \Cue{G} &
Two-sample Welch \textit{t}-test &
Two independent groups are compared and equal variances should not be assumed. &
$H_0$: $\mu_1 = \mu_2$. &
\Rc{t.test(y \textasciitilde{} g, data = df)} &
Usually the default. Assumes independence and approximate within-group normality. \\

7 & \Cue{G} &
Two-sample Student \textit{t}-test &
Two independent groups are compared and equal variances are defensible. &
$H_0$: $\mu_1 = \mu_2$. &
\Rc{t.test(y \textasciitilde{} g, data = df,}\newline
\Rc{\ \ var.equal = TRUE)} &
Use only when group spreads are sufficiently similar. \\

7 & \Cue{G} &
Paired \textit{t}-test &
Two linked measurements are compared within the same sampling units. &
$H_0$: mean paired difference $= 0$. &
\Rc{t.test(x, y, paired = TRUE)} &
Assumptions apply to within-pair differences, not the raw before/after values. \\

7 & \Cue{Rk} &
Wilcoxon signed-rank test &
One-sample or paired problem when a rank-based alternative is more defensible. &
$H_0$: the location shift is 0. &
\Rc{wilcox.test(x, mu = m0)}\newline
\Rc{wilcox.test(x, y, paired = TRUE)} &
Use for one-sample and paired designs. Still requires the correct design logic. \\

7 & \Cue{Rk} &
Wilcoxon rank-sum / Mann--Whitney test &
Two independent groups are compared without relying on Gaussian assumptions. &
$H_0$: the two groups come from the same distributional location. &
\Rc{wilcox.test(y \textasciitilde{} g, data = df)} &
Rank-based alternative to the two-sample \textit{t}-test. \\

\rowcolor{groupbg}\multicolumn{7}{l}{\textbf{ANOVA family and rank-based group comparisons}} \\
8 & \Cue{G} &
One-way ANOVA &
One continuous response is compared across more than two independent groups. &
$H_0$: \(\mu_1 = \mu_2 = \cdots = \mu_k\). &
\Rc{aov(y \textasciitilde{} g, data = df)} &
Check residual normality and equal variance. This is the multi-group mean-comparison model. \\

8 & \Cue{G} &
Factorial / two-way ANOVA &
Two or more categorical predictors are fitted together without continuous covariates. &
$H_0$: no main effects and, if included, no interaction effect. &
\Rc{aov(y \textasciitilde{} A * B, data = df)} &
Use \Rc{A + B} for additive structure and \Rc{A * B} when the interaction is part of the question. \\

8,14 & \Cue{G} &
ANCOVA as a linear model &
One continuous response is modelled using both categorical and continuous predictors. &
Model target: adjusted group differences plus covariate slope(s). &
\Rc{lm(y \textasciitilde{} group + x, data = df)} &
ANCOVA is a standard linear model with one factor and one or more covariates. \\

8 & \Cue{G} &
Tukey's HSD post hoc comparison &
One-way ANOVA is significant and pairwise mean differences among groups are needed. &
$H_0$: each pairwise group mean difference $= 0$. &
\Rc{TukeyHSD(fit)} &
Use after a fitted \Rc{aov()} object. Controls family-wise error for all pairwise comparisons. \\

8 & \Cue{Rk} &
Kruskal--Wallis test &
More than two independent groups are compared with a rank-based omnibus test. &
$H_0$: groups have the same distributional location. &
\Rc{kruskal.test(y \textasciitilde{} g, data = df)} &
Rank-based alternative to one-way ANOVA. \\

8 & \Cue{Rk} &
Pairwise Wilcoxon comparisons &
Kruskal--Wallis or multi-group rank comparisons require pairwise follow-up. &
$H_0$: each pair of groups has the same location. &
\Rc{pairwise.wilcox.test(y, g,}\newline
\Rc{\ \ p.adjust.method = "BH")} &
Used as a follow-up rather than as the primary omnibus test. \\

\rowcolor{headbg}\multicolumn{7}{l}{\color{headfg}\textbf{CONTINUOUS MODELING}} \\
\rowcolor{groupbg}\multicolumn{7}{l}{\textbf{Association and correlation}} \\
9 & \Cue{Asc} &
Pearson correlation &
Two continuous variables are paired and a linear association is of interest without fitting a response model. &
$H_0$: $\rho = 0$. &
\Rc{cor.test(x, y, method = "pearson")} &
Appropriate for approximately linear continuous relationships. \\

9 & \Cue{Asc} &
Spearman rank correlation &
Two paired variables show a monotonic but not necessarily linear association, or ranks are more appropriate. &
$H_0$: $\rho_s = 0$. &
\Rc{cor.test(x, y, method = "spearman")} &
Works for ordinal or non-Gaussian paired data. \\

9 & \Cue{Asc} &
Kendall's $\tau$ &
Two paired variables are associated, often with small samples or many ties. &
$H_0$: $\tau = 0$. &
\Rc{cor.test(x, y, method = "kendall")} &
Often more conservative than Spearman. \\

\rowcolor{groupbg}\multicolumn{7}{l}{\textbf{Gaussian regression: mean structure}} \\
\rowcolor{diagbg}11 & \Cue{G} &
Residual diagnostics for fitted Gaussian models &
A regression model has been fitted and the adequacy of mean and variance structure must be checked. &
No null; the aim is to inspect residual structure. &
\Rc{plot(mod, which = 1)}\newline
\Rc{plot(mod, which = 2)} &
Residual-vs-fitted checks mean structure and spread. Q--Q checks residual normality. Independence comes from design. \\

12 & \Cue{G} &
Simple linear regression &
One continuous response is modelled as a linear function of one continuous predictor. &
Model target: mean response changes linearly with the predictor. &
\Rc{lm(y \textasciitilde{} x, data = df)} &
Inference on slopes uses the fitted linear model, not a correlation test in disguise. \\

\rowcolor{flowbg}12 & \Cue{G} &
Confidence and prediction intervals &
A fitted linear model is used for mean-response uncertainty or individual-response prediction. &
No null; the aim is interval estimation. &
\Rc{predict(mod, newdata = nd,}\newline
\Rc{\ \ interval = "confidence")}\newline
\Rc{predict(mod, newdata = nd,}\newline
\Rc{\ \ interval = "prediction")} &
Confidence intervals describe the mean line; prediction intervals describe future observations and are wider. \\

\rowcolor{diagbg}12,13 & \Cue{G} &
Breusch--Pagan test &
Heteroscedasticity in a fitted linear model needs a formal supporting check. &
$H_0$: residual variance is constant. &
\Rc{lmtest::bptest(mod)} &
Use as supporting evidence after reading residual plots. \\

\rowcolor{diagbg}12,17 & \Cue{G} &
Influence diagnostics &
Single observations may be driving the fitted model disproportionately. &
No single null; the aim is to identify leverage and influence. &
\Rc{cooks.distance(mod)}\newline
\Rc{dffits(mod)}\newline
\Rc{plot(mod, which = 5)} &
Inspect Cook's distance, DFFITS, leverage, and residual-vs-leverage plots before acting on a point. \\

13 & \Cue{G} &
Polynomial regression &
A linear model is too rigid but the problem can still be handled within the \Rc{lm()} framework. &
Model target: mean structure includes curvature. &
\Rc{lm(y \textasciitilde{} x + I(x\^{}2), data = df)}\newline
\Rc{lm(y \textasciitilde{} poly(x, degree = k), data = df)} &
Still linear in the coefficients. Compare nested degrees by \Rc{anova()} or \Rc{AIC()}. \\

14 & \Cue{G} &
Multiple regression &
One continuous response depends on several predictors fitted simultaneously. &
Model target: partial effects after adjusting for the other predictors. &
\Rc{lm(y \textasciitilde{} x1 + x2 + x3, data = df)} &
Interpret coefficients conditionally. Model specification matters as much as fit. \\

14,15 & \Cue{G} &
Interaction model &
The effect of one predictor may depend on another predictor. &
$H_0$: the interaction coefficient $= 0$. &
\Rc{lm(y \textasciitilde{} x1 * x2, data = df)} &
Use \Rc{*} to include both main effects and the interaction. Respect hierarchy. \\

\shortstack[l]{14,15,\\17} & \Cue{G} &
Nested model comparison for Gaussian linear models &
A simpler and a more complex \Rc{lm()} model are nested and need formal comparison. &
$H_0$: added term(s) do not reduce residual variance beyond chance. &
\Rc{anova(mod\_small, mod\_large,}\newline
\Rc{\ \ test = "F")} &
Appropriate when one model is a special case of the other. \\

\rowcolor{headbg}\multicolumn{7}{l}{\color{headfg}\textbf{MODEL SELECTION AND MODERN WORKFLOWS}} \\
\rowcolor{flowbg}\shortstack[l]{14,17,\\24} & \Cue{Any} &
Akaike Information Criterion &
Several biologically plausible models need relative comparison without a single null hypothesis. &
Lower AIC indicates better trade-off between fit and complexity. &
\Rc{AIC(mod1, mod2, mod3)} &
Useful for non-nested and nested candidate models. Compare only models fitted to the same rows, response data, and response scale. \\

\rowcolor{diagbg}16 & \Cue{G} &
Variance Inflation Factor (VIF) &
Predictors in a multiple regression may be collinear. &
No null; the aim is to quantify variance inflation from shared predictor information. &
\Rc{car::vif(mod)} &
Large VIFs indicate unstable coefficient interpretation. Use alongside correlation matrices and scatterplots. \\

\rowcolor{groupbg}\multicolumn{7}{l}{\textbf{Dependence and grouped structure}} \\
\rowcolor{flowbg}18 & \Cue{Grp} &
Design check for pseudoreplication &
Measurements may have been treated as independent when the design does not support that assumption. &
No formal residual null; the target is correct unit-of-analysis logic. &
\Rc{No single function} &
This is a design problem, not a residual issue. Mixed models handle structured dependence, but they do not rescue true pseudoreplication. \\

19 & \Cue{Grp} &
Linear mixed model &
Gaussian response with grouped, repeated, or clustered observations. &
Model target: fixed effects plus random-group variation. &
\Rc{lme4::lmer(y \textasciitilde{} x + (1 | group),}\newline
\Rc{\ \ data = df)} &
Use when dependence is structured by design, for example shared sites, individuals, years, or plots. This is not a repair for \(n=1\) pseudoreplication. \\

19 & \Cue{Grp} &
Generalised linear mixed model &
Non-Gaussian response with grouped, repeated, or clustered observations. &
Model target: non-Gaussian mean structure plus random effects. &
\Rc{lme4::glmer(y \textasciitilde{} x + (1 | group),}\newline
\Rc{\ \ family = ..., data = df)} &
Extends mixed-model logic to binomial or count responses. \\

\rowcolor{diagbg}19 & \Cue{Grp} &
Intra-class correlation (ICC) &
Variance partitioning is needed after fitting a mixed model. &
No null; the target is the proportion of variance attributable to groups. &
\Rc{VarCorr(mod)}, then compute ICC &
Used to summarise the strength of dependence at the grouping level. \\

\rowcolor{groupbg}\multicolumn{7}{l}{\textbf{Generalised linear models and proportion/count responses}} \\
20 & \Cue{Bin} &
One-sample proportion test &
A single observed proportion is compared with a null expectation. &
$H_0$: $p = p_0$. &
\Rc{prop.test(x = s, n = n, p = p0)} &
Presented as a bridge to binomial GLMs when predictors are absent. \\

20 & \Cue{Bin} &
Two-sample proportion test &
Two binomial proportions are compared without additional predictors. &
$H_0$: $p_1 = p_2$. &
\Rc{prop.test(x = c(s1, s2), n = c(n1, n2))}\newline
\Rc{prop.test(matrix\_counts)} &
Useful for simple comparisons; move to a binomial GLM when predictors enter. \\

20 & \Cue{Bin} &
Logistic regression (Bernoulli response) &
Binary presence/absence or success/failure response with predictors. &
Model target: \(\mathrm{logit}\{P(y=1)\} = \beta_0 + \beta_1x + \cdots\). &
\Rc{glm(y \textasciitilde{} x1 + x2,}\newline
\Rc{\ \ family = binomial, data = df)} &
Use \Rc{type = "response"} in \Rc{predict()} to recover fitted probabilities. \\

20 & \Cue{Bin} &
Binomial GLM for grouped counts &
Successes and failures are aggregated for each row. &
Model target: \(\mathrm{logit}(p) = \beta_0 + \beta_1x + \cdots\). &
\Rc{glm(cbind(success, failure) \textasciitilde{} x,}\newline
\Rc{\ \ family = binomial, data = df)} &
Same link and interpretation as logistic regression, but with grouped trials. \\

20 & \Cue{Cnt} &
Poisson regression &
Count response with mean approximately equal to the variance. &
Model target: \(\log\{E[y]\} = \beta_0 + \beta_1x + \cdots\). &
\Rc{glm(count \textasciitilde{} x,}\newline
\Rc{\ \ family = poisson, data = df)} &
Check for overdispersion before trusting Poisson standard errors. \\

20 & \Cue{Cnt} &
Quasi-Poisson regression &
Count response shows overdispersion relative to Poisson. &
Model target: \(\log\{E[y]\} = \beta_0 + \beta_1x + \cdots\), with overdispersed variance. &
\Rc{glm(count \textasciitilde{} x,}\newline
\Rc{\ \ family = quasipoisson, data = df)} &
Retains the log-link mean model but adjusts the variance. \\

\rowcolor{groupbg}\multicolumn{7}{l}{\textbf{Flexible and nonlinear regression}} \\
21 & \Cue{Sm} &
Gaussian GAM &
Relationship is smooth and nonlinear, but a fixed polynomial form is not appropriate. &
Model target: smooth mean function estimated from the data. &
\Rc{mgcv::gam(y \textasciitilde{} s(x),}\newline
\Rc{\ \ data = df, method = "REML")} &
Use \Rc{mgcv::gam()} or load \Rc{library(mgcv)}. Use cyclic bases such as \Rc{bs = "cc"} for seasonal smooths. \\

22 & \Cue{NL} &
Nonlinear least-squares regression &
Mean structure follows a named nonlinear function with interpretable parameters. &
Model target: explicit nonlinear response curve, for example Michaelis--Menten. &
\Rc{nls(y \textasciitilde{} f(x, pars),}\newline
\Rc{\ \ start = list(...), data = df)} &
Requires sensible starting values and convergence checking. \\

\rowcolor{flowbg}22 & \Cue{NL} &
Linear-vs-nonlinear model comparison &
A nonlinear model is compared with a simpler linear alternative. &
Relative support for alternative mean structures. &
\Rc{anova(lm\_mod, nls\_mod)}\newline
\Rc{AIC(lm\_mod, nls\_mod)} &
Interpret with care because \Rc{nls()} and \Rc{lm()} may reflect different biological questions. \\

23 & \Cue{Q} &
Quantile regression &
Predictor effects may differ across the response distribution rather than only at the mean. &
Model target: conditional quantile(s), not conditional mean alone. &
\Rc{quantreg::rq(y \textasciitilde{} x,}\newline
\Rc{\ \ tau = 0.5, data = df)} &
Typical choices are \Rc{tau = 0.1}, \Rc{0.5}, and \Rc{0.9}. \\

23 & \Cue{Q} &
Quantile-slope comparison &
Quantile regressions have been fitted at several \(\tau\) values and slope differences need testing. &
$H_0$: slopes do not differ across the fitted quantiles. &
\Rc{anova(rq\_fit\_multi\_tau)} &
Useful when upper, median, and lower response boundaries behave differently. \\

\rowcolor{groupbg}\multicolumn{7}{l}{\textbf{Prediction, evaluation, and regularisation}} \\
\rowcolor{flowbg}17,24 & \Cue{Pred} &
Hold-out RMSE &
Prediction quality is assessed on data not used for fitting. &
No null; lower RMSE indicates better predictive performance. &
\Rc{pred <- predict(mod, newdata = test)}\newline
\Rc{sqrt(mean((obs - pred)\^{}2))} &
Prediction and explanation are different goals; a more interpretable model need not minimise prediction error. \\

\rowcolor{flowbg}17 & \Cue{Pred} &
K-fold cross-validation for Gaussian models &
Generalisation error is estimated by repeated train/test splits across folds. &
No null; compare average out-of-fold RMSE across candidate models. &
fit with \Rc{lm()} inside each fold,\newline
use \Rc{predict()}, then summarise RMSE &
Implemented manually in the chapter rather than through a single package wrapper. \\

25 & \Cue{Pred} &
Ridge regression &
Many correlated predictors are present and coefficient shrinkage is acceptable. &
Model target: stable prediction with \(L_2\) penalty. &
\Rc{glmnet::cv.glmnet(X, y, alpha = 0)} &
Shrinks coefficients toward zero but typically keeps all of them non-zero. \\

25 & \Cue{Pred} &
Lasso regression &
Prediction and variable selection are both of interest. &
Model target: sparse regression with \(L_1\) penalty. &
\Rc{glmnet::cv.glmnet(X, y, alpha = 1)} &
Can set some coefficients exactly to zero. \\

25 & \Cue{Pred} &
Elastic net &
Predictors are correlated and a compromise between ridge and lasso is needed. &
Model target: mixed \(L_1/L_2\) penalty. &
\Rc{glmnet::cv.glmnet(X, y,}\newline
\Rc{\ \ alpha = 0.5)} &
Choose \Rc{alpha} between 0 and 1. Particularly useful when predictors cluster strongly. \\

\end{longtable}
}

\vspace{4pt}
\noindent\hspace*{\tableoffset}%
\begin{minipage}{\tableblockwidth}
\rule{\tableblockwidth}{0.5pt}
\end{minipage}
\vspace{3pt}

\noindent\hspace*{\tableoffset}%
\begin{minipage}{\tableblockwidth}
\begin{multicols}{3}
\footnotesize

\noindent\textbf{Design first.} Choose the method from the data structure and not from the software options. The same response variable may require a \textit{t}-test, ANOVA, regression, mixed model, or GLM depending on independence, grouping, pairing, and whether predictors are categorical or continuous.

\medskip
\noindent\textbf{Response scale.} Gaussian methods such as \Rc{t.test()}, \Rc{aov()}, and \Rc{lm()} assume that the analysis is defensible on the chosen response scale. For \Rc{aov()} and \Rc{lm()}, this is checked mainly through residuals. For \Rc{t.test()}, checks apply to the relevant raw values or paired differences. If the assumptions fail, transform the response or move to a method better matched to the data.

\medskip
\noindent\textbf{Model hierarchy.} When an interaction is included, keep its main effects. When grouped structure is present, fit the grouping rather than pretending the rows are independent.

\medskip
\noindent\textbf{Model comparison.} Use nested \Rc{anova(..., test = "F")} when one model is a special case of another. Use \Rc{AIC()} for broader candidate-model comparison, but not across different response scales such as \(y\) versus \(\log(y)\) or across models fitted to different numbers of observations after NA-driven row loss. Use hold-out or cross-validated RMSE when prediction is the objective.

\noindent\textbf{Cue legend.} \Cue{G} Gaussian continuous; \Cue{Rk} rank-based; \Cue{Asc} association; \Cue{Bin} binary or proportion; \Cue{Cnt} count; \Cue{Grp} grouped or repeated; \Cue{Sm} smooth model; \Cue{NL} nonlinear; \Cue{Q} quantile; \Cue{Pred} predictive workflow; \Cue{Any} model-comparison workflow.

\medskip
\noindent\textbf{Icons.} Icons are mnemonics: \CueIcon{G} mean-based Gaussian, \CueIcon{Rk} rank order, \CueIcon{Asc} paired association, \CueIcon{Bin} binary response, \CueIcon{Cnt} counts, \CueIcon{Grp} clustered observations, \CueIcon{Sm} smooth function, \CueIcon{NL} explicit nonlinear form, \CueIcon{Q} quantiles, \CueIcon{Pred} prediction, \CueIcon{Any} model comparison.

\medskip
\noindent\textbf{Row tinting.} Cream rows mark diagnostics or assumption checks. Green rows mark workflow, comparison, or prediction entries rather than primary fitted procedures.

\medskip
\noindent\textbf{Abbreviations.} GLM = generalised linear model; GAM = generalised additive model; ICC = intra-class correlation; RMSE = root mean squared error; VIF = variance inflation factor.

\medskip
\noindent\textbf{Packages beyond base R used in the course.}
\begin{itemize}[leftmargin=*, itemsep=1pt, topsep=2pt]
\item \Rc{car} for \Rc{leveneTest()} and \Rc{vif()}
\item \Rc{lmtest} for \Rc{bptest()}
\item \Rc{lme4} for \Rc{lmer()} and \Rc{glmer()}
\item \Rc{mgcv} for \Rc{gam()}
\item \Rc{quantreg} for \Rc{rq()}
\item \Rc{glmnet} for ridge, lasso, and elastic net
\end{itemize}

\medskip
\noindent\textbf{Not every chapter adds a new test.} Chapters 1--5, 10, 18, 24, and 26 mainly refine framing, design logic, model choice, interpretation, prediction, and workflow rather than introducing a brand-new inferential procedure.

\medskip
\noindent\textbf{Course synthesis companion.} The chapter summary is in \Rc{cheatsheet-course-synthesis.pdf}, leaving this sheet procedural and faster to scan.

\end{multicols}
\end{minipage}

\vspace{5pt}
\noindent\hspace*{\tableoffset}%
\begin{minipage}{\tableblockwidth}
\rule{\tableblockwidth}{0.5pt}
\end{minipage}
\vspace{4pt}

\noindent\hspace*{\tableoffset}%
\begin{minipage}{\tableblockwidth}
\noindent{\large\bfseries Safer Use Checks}

\vspace{3pt}

\begin{tcbraster}[
  raster columns=3,
  raster equal height=rows,
  raster column skip=4mm,
  raster left skip=0pt,
  raster right skip=0pt
]
\begin{infopanelbox}{diagbg}{Sampling-Unit Check}
\begin{itemize}[leftmargin=*, itemsep=2pt, topsep=2pt]
\item The sampling unit is the smallest independently sampled biological or experimental replicate to which inference is intended to apply.
\item One-sample procedures: one value per independent sampling unit.
\item Two-sample procedures: one row per independent sampling unit, and no sampling unit appears in both groups.
\item Paired procedures: the sampling unit is the matched pair; analyse within-pair differences, not raw columns as independent samples.
\item ANOVA and ANCOVA: groups must contain independent sampling units unless grouping is modelled explicitly.
\item Correlation: both variables are measured on the same sampling units; neither variable is a replicated treatment label in disguise.
\item Regression: repeated rows from the same site, tank, plot, year, or individual are not independent sampling units.
\item Mixed and grouped models: use when rows are clustered within higher-level sampling units. They partition variance and dependence; they do not create new independent replicates.
\end{itemize}
\end{infopanelbox}
\begin{warnpanelbox}{Common Misuse Warnings}
\begin{itemize}[leftmargin=*, itemsep=2pt, topsep=2pt]
\item Do not choose a test from the software menu before identifying the sampling unit and dependence structure.
\item Do not treat technical replicates, subsamples, or repeated measures as biological replicates.
\item Do not use Student's two-sample \textit{t}-test by default; Welch is usually the safer default.
\item Do not read a significant omnibus ANOVA as evidence that every pair of groups differs.
\item Do not read a correlation coefficient as a causal effect or as a substitute for a regression with a defined response.
\item Do not use residual normality tests to rescue a misspecified mean structure; fix the model first.
\item Do not drop main effects when an interaction is retained.
\item Do not interpret collinear coefficients as stable biological effects without checking VIFs, plots, and study design.
\item Do not use AIC across different response datasets or silently different missing-data subsets.
\item Do not treat lower predictive error as proof of better biological explanation.
\end{itemize}
\end{warnpanelbox}
\begin{infopanelbox}{reportbg}{Minimum Reporting Checklist}
\begin{itemize}[leftmargin=*, itemsep=2pt, topsep=2pt]
\item \textit{t}-tests and ANOVA: group summaries, effect direction, test statistic, degrees of freedom, \(p\)-value, and interval where relevant.
\item Rank tests: sample summaries, test statistic, \(p\)-value, and a clear statement that inference is rank-based.
\item Correlation: coefficient type, coefficient estimate, sample size, confidence interval if available, and \(p\)-value.
\item Linear models: formula, coefficients, standard errors, test statistics, \(p\)-values, uncertainty intervals, and residual-diagnostic conclusion.
\item Mixed models: fixed-effect estimates, random-effect structure, variance components or ICC, and the grouping variable.
\item GLMs: family and link, coefficients on the model scale, transformed interpretation where needed, and any dispersion check.
\item GAMs and nonlinear models: fitted form, smoothing or parameter details, uncertainty summary, and graphical diagnostic conclusion.
\item Predictive comparisons: data split or resampling scheme, performance metric, and whether the goal is explanation or prediction.
\end{itemize}
\end{infopanelbox}
\end{tcbraster}
\end{minipage}

\end{document}
